% ==============================================================
%  Chapter 2 -- Theoretical Framework
% ==============================================================

\chapter{Theoretical Framework}
\label{ch:framework}

This chapter introduces the theoretical background on which the rest of
the thesis relies. We start with Decision Support Systems and their graphical formulation as Influence Diagrams, fix the formal
definitions that are used in subsequent chapters and motivate the
parameter-estimation problem that we tackle. We then present
Estimation of Distribution Algorithms, justify our design choices, including the specific
variants and the convergence guarantees that ground the empirical study of Chapter~4.

% ==============================================================
\section{Decision Support Systems and Influence Diagrams}
% ==============================================================

A \emph{Decision Support System} (DSS) is a model or piece of
software designed to assist decision-makers by facilitating access to
relevant data and by applying analytical techniques to
poorly structured problems~\cite{insua2005,shenoy2009eolss}.
The role of a DSS is not to replace the decision-maker but to provide
structured access to information and a systematic way of evaluating
alternatives before a final decision is reached.

Among the many graphical formalisms available, \emph{Influence
Diagrams} (IDs) have established themselves as one of the most
compact and expressive
representations~\cite{howard2005influence,koller2009pgm}. An ID
is a directed acyclic graph (DAG) whose nodes encode the three fundamental
ingredients of a sequential decision problem: random variables,
controllable choices and preferences, and whose arcs encode
conditional dependencies and informational precedence.

\subsection{Formal definition}

\begin{definition}[Influence Diagram]
\label{def:id}
An \emph{Influence Diagram} is a tuple
$\mathcal{D} = (\mathbf{C}, \mathbf{D}, V, \mathcal{E})$ where:
\begin{itemize}
    \item $\mathbf{C} = \{C_1,\ldots,C_n\}$ is a finite set of
          \emph{chance nodes}, each modelling a discrete random
          variable with finite domain $\Omega_{C_i}$;
    \item $\mathbf{D} = \{D_1,\ldots,D_k\}$ is a finite set of
          \emph{decision nodes}, each with a finite action set
          $\mathcal{A}_{D_i}$;
    \item $V$ is a single \emph{utility} (value) node encoding the
          preferences of the decision-maker;
    \item $\mathcal{E}$ is a set of directed arcs split into
          probabilistic, informational and functional arcs.
\end{itemize}
\end{definition}

For the ID to admit a well-defined evaluation procedure the graph
must be \emph{regular}. That is, a DAG where the value node has no successors and the decisions can be linearly ordered by a directed path that traverses all of them. Throughout this work all IDs are
assumed to be regular.

\subsubsection*{A canonical example: the oil wildcatter}
\label{sec:oil-wildcatter}

To make the formalism concrete we adopt the canonical \emph{oil
wildcatter} problem of Howard and Matheson~\cite{howard2005influence},
also reproduced in the review of Bielza, G{\'o}mez and
Shenoy~\cite{bielza2011review}. A prospector must decide whether to
drill at a given site. Before drilling, she may pay for a seismic
test whose result is informative about the underlying amount of oil.
The graphical structure is shown in Figure~\ref{fig:oilid}.

\begin{figure}[H]
    \centering
    \begin{tikzpicture}[
        >=Stealth, node distance=14mm and 18mm, font=\small,
        chance/.style={draw, ellipse, minimum width=18mm, minimum height=8mm},
        decision/.style={draw, rectangle, minimum width=18mm, minimum height=8mm},
        utility/.style={draw, diamond, aspect=2, minimum width=14mm,
                        inner sep=2pt},
        arr/.style={->, thick},
        infoarr/.style={->, thick, dashed}
    ]
        \node[decision] (T)              {Test?};
        \node[chance,   right=of T]  (S) {Seismic};
        \node[decision, right=of S]  (D) {Drill?};
        \node[chance,   above=of S]  (O) {Oil};
        \node[utility,  below=10mm of S] (V) {Profit};

        \draw[arr] (T) -- (S);
        \draw[arr] (O) -- (S);
        \draw[infoarr] (T) -- (D);
        \draw[infoarr] (S) -- (D);
        \draw[arr] (T) -- (V);
        \draw[arr] (D) -- (V);
        \draw[arr] (O.east) to[bend left=20] (V.east);
    \end{tikzpicture}
    \caption{Influence diagram for the oil-wildcatter problem.
    Ellipses are chance nodes, rectangles are decision nodes, the
    diamond is the utility node. Solid arrows denote probabilistic
    or functional dependencies, dashed arrows denote informational
    precedence (what is known when each decision is taken).}
    \label{fig:oilid}
\end{figure}

The qualitative structure is read off the picture: the seismic test
result depends on whether the test was performed and on the true
amount of oil; the drilling decision is taken after observing both;
the profit depends on the drilling action, on the true amount of oil
(success or dry well) and on whether the seismic test was paid for.
The next subsection puts numbers to this skeleton.

\subsection{Quantitative parameterisation}

Each chance node $C_i$ carries a \emph{conditional probability
table} (CPT) $\theta_{C_i}$ that assigns to every configuration of
its probabilistic parents $\mathrm{Pa}(C_i)$ a probability distribution over
$\Omega_{C_i}$:
\[
   \theta_{C_i}\bigl(\cdot \mid \mathbf{pa}\bigr) \;\in\;
   \Delta(\Omega_{C_i})
   \qquad
   \forall \mathbf{pa}\in\prod_{C_j\in\mathrm{Pa}(C_i)}\Omega_{C_j},
\]
where $\Delta(\Omega_{C_i})$ denotes the probability simplex over the
states of $C_i$. The utility node carries a real-valued mapping
\[
   u : \prod_{X\in\mathrm{Pa}(V)}\Omega_X \longrightarrow \mathbb{R},
\]
where $\Omega_{D_i}\equiv\mathcal{A}_{D_i}$ for any decision parent.
The function $u$ encodes the preferences of the decision-maker in
the sense of Multi-Attribute Utility Theory, and when its arguments
satisfy the usual independence conditions it admits an additive or
multiplicative decomposition that drastically simplifies its
elicitation~\cite{riosinsua2002maut}. The collection of all CPTs
and the utility table is the \emph{parameter vector} of the
diagram, written $\boldsymbol{\Theta}$.

\begin{remark}[Combinatorial explosion]
\label{rem:explosion}
The number of free entries of a CPT grows multiplicatively with the
parents: a node with $m$ parents, each taking $r$ states, and $r$
own states contributes $r^{m}(r-1)$ free parameters. This
exponential growth is the source of the elicitation bottleneck
discussed in Chapter~\ref{ch:introduction} and the reason why
parameter optimisation rather than exhaustive elicitation is
necessary.
\end{remark}

A \emph{decision rule} for $D_i$ is a mapping
$\delta_i : \prod_{C_j\in\mathrm{Pa}(D_i)}\Omega_{C_j}
\longrightarrow \mathcal{A}_{D_i}$
that prescribes an action for every observable context. The
collection $\boldsymbol{\delta}=(\delta_1,\ldots,\delta_k)$ over all
decisions is called a \emph{policy}. The expected utility of a
policy under a parameter vector $\boldsymbol{\Theta}$ is
\begin{equation}
\label{eq:eu}
   \mathrm{EU}(\boldsymbol{\delta}\mid\boldsymbol{\Theta}) =
   \sum_{\mathbf{c}} u\!\bigl(\mathbf{c},
                              \boldsymbol{\delta}(\mathbf{c})\bigr)
   \prod_{C_i\in\mathbf{C}}
   \theta_{C_i}\!\bigl(c_i \mid \mathbf{pa}(c_i)\bigr),
\end{equation}
where $\mathbf{c}=(c_1,\ldots,c_n)$ is a complete \emph{scenario}. For each scenario,
$\mathbf{pa}(c_i)$ denotes the values its parents take within that same
scenario and $\boldsymbol{\delta}(\mathbf{c})$ the actions prescribed
by the policy. Thus each scenario is weighted by its joint probability
(the product of CPTs) and scored by its utility, and the expectation
sums these contributions over every possible state of the space.
The \emph{optimal policy} is
$\boldsymbol{\delta}^{\!*}(\boldsymbol{\Theta}) =
\arg\max_{\boldsymbol{\delta}}
\mathrm{EU}(\boldsymbol{\delta}\mid\boldsymbol{\Theta})$.

\noindent The classical algorithm of
Shachter~\cite{shachter1986evaluating} computes
$\boldsymbol{\delta}^{\!*}(\boldsymbol{\Theta})$ in polynomial time
for any fixed $\boldsymbol{\Theta}$ by successively reversing arcs
and eliminating nodes from the graph.

\subsubsection*{What the parameters and their evaluation look like}
\label{sec:oil-tables}

To make the three pieces of $\boldsymbol{\Theta}$ concrete we add
illustrative numbers to the oil-wildcatter problem of
Figure~\ref{fig:oilid}. The chance node \textsf{Oil} takes three
states (\textsc{dry}, \textsc{wet}, \textsc{soaking}); the chance
node \textsf{Seismic} takes three states (\textsc{none},
\textsc{open}, \textsc{closed}) and depends on the true \textsf{Oil}
amount and on whether the test was performed. The CPT of
\textsf{Seismic} conditioned on $\mathsf{Test}=\textsc{yes}$ is shown
in Table~\ref{tab:oil-cpt}; we omit the trivial column
$\mathsf{Test}=\textsc{no}$, in which the test returns a fixed
\textsc{none}.

\begin{table}[H]
\centering\small
\begin{tabular}{l|ccc}
\hline
$P(\mathsf{Seismic} \mid \mathsf{Oil}, \mathsf{Test}=\textsc{yes})$
              & \textsc{none} & \textsc{open} & \textsc{closed} \\
\hline
$\mathsf{Oil}=\textsc{dry}$     & 0.60 & 0.30 & 0.10 \\
$\mathsf{Oil}=\textsc{wet}$     & 0.30 & 0.40 & 0.30 \\
$\mathsf{Oil}=\textsc{soaking}$ & 0.10 & 0.30 & 0.60 \\
\hline
\end{tabular}
\caption{Conditional probability table of the chance node
\textsf{Seismic} given its parents. Each row is a probability
distribution over the three possible test outcomes.}
\label{tab:oil-cpt}
\end{table}

The utility node \textsf{Profit} depends on \textsf{Drill},
\textsf{Oil} and \textsf{Test}; Table~\ref{tab:oil-util} lists its
values on an arbitrary monetary scale (in our convention, the
utility lies in $[0,10]$).

\begin{table}[H]
\centering\small
\begin{tabular}{lcc|cc}
\hline
& & & $\mathsf{Test}=\textsc{no}$ & $\mathsf{Test}=\textsc{yes}$\\
\hline
$\mathsf{Drill}=\textsc{no}$  & & $\mathsf{Oil}=$ any       & 5.0 & 4.0 \\
\hline
$\mathsf{Drill}=\textsc{yes}$ & & $\mathsf{Oil}=\textsc{dry}$     & 0.0 & 0.0 \\
                              & & $\mathsf{Oil}=\textsc{wet}$     & 6.0 & 5.0 \\
                              & & $\mathsf{Oil}=\textsc{soaking}$ & 10.0& 9.0 \\
\hline
\end{tabular}
\caption{Utility table of \textsf{Profit} on the $[0,10]$ scale.
Drilling a dry well loses money; drilling a soaking well is the most
profitable outcome; the test reduces the utility by a constant cost
(here, one unit).}
\label{tab:oil-util}
\end{table}

Given these two ingredients and a prior on \textsf{Oil}
(say $P(\textsc{dry}) = 0.5,\, P(\textsc{wet}) = 0.3,\,
P(\textsc{soaking}) = 0.2$), running the Shachter algorithm produces
a third table: the expected utility of each action of $\mathsf{Drill}$
\emph{conditioned on the information available at that point} (the
seismic result, if a test was paid for). Table~\ref{tab:oil-eu} shows
this evaluation output. The action of larger expected utility in each
row is the optimal one for that information state, which is what
makes the table the carrier of the \emph{decision rules} of the model.

\begin{table}[H]
\centering\small
\begin{tabular}{l|cc}
\hline
$\mathsf{EU}\bigl(\mathsf{Drill}\mid \mathsf{Test}, \mathsf{Seismic}\bigr)$
                                 & $\mathsf{Drill}=\textsc{no}$ & $\mathsf{Drill}=\textsc{yes}$\\
\hline
$\mathsf{Test}=\textsc{no}$      & 5.00 & \textbf{4.40}\\[2pt]
\hline
$\mathsf{Test}=\textsc{yes},\;\mathsf{Seismic}=\textsc{none}$
                                 & 4.00 & \textbf{2.36}\\
$\mathsf{Test}=\textsc{yes},\;\mathsf{Seismic}=\textsc{open}$
                                 & 4.00 & \textbf{4.94}\\
$\mathsf{Test}=\textsc{yes},\;\mathsf{Seismic}=\textsc{closed}$
                                 & 4.00 & \textbf{6.65}\\
\hline
\end{tabular}
\caption{Expected utilities produced by the Shachter evaluation of
the oil-wildcatter ID under the parameters of
Tables~\ref{tab:oil-cpt}--\ref{tab:oil-util}.
The values in bold are the optimal action of $\mathsf{Drill}$ for
that information state. The whole rightmost column is what a domain
expert would summarise as ``decision rules'': drill if and only if
the seismic test is open or closed, otherwise abstain.}
\label{tab:oil-eu}
\end{table}

The point of this thesis is precisely the \emph{inverse} of this
procedure. The expert does \emph{not} provide
Tables~\ref{tab:oil-cpt} and~\ref{tab:oil-util}: instead, she only
supplies a few entries of Table~\ref{tab:oil-eu} as decision rules
(``drill if the seismic is closed'', ``don't drill if the seismic is
none''). The optimiser must recover a pair of tables
$(\theta_{C_i}, u)$ \emph{compatible} with those rules, in the sense
that re-running Shachter on the recovered parameters reproduces the
expert's choices.

% ==============================================================
\section{The Knowledge-Acquisition Problem}
% ==============================================================

In practice the structure of the ID is elicited from the expert in a
small number of modelling sessions, but the numerical
parameterisation $\boldsymbol{\Theta}$ is far harder to obtain. The
classical cycle of~\cite{bielza2000issues,bielza2010modeling} proceeds in
five stages — problem identification, enumeration of alternatives,
modelling, evaluation and sensitivity analysis — and the modelling
stage alone accounts for the bulk of the cost. As anticipated in
Remark~\ref{rem:explosion}, the size of $\boldsymbol{\Theta}$ grows
exponentially with the connectivity of the diagram, so an exhaustive
elicitation is unrealistic for any diagram of practical size.

We assume instead that the expert can provide a \emph{small set of
decision rules}.

\begin{definition}[Expert rule base]
\label{def:rules}
A \emph{rule base} is a finite set
\[
   \mathcal{R} = \bigl\{(\mathbf{x}^{(l)}, D_{i_l}, a^{(l)})\bigr\}_{l=1}^{L},
\]
where $\mathbf{x}^{(l)}$ is a (possibly partial) observation vector
on the informational parents of decision $D_{i_l}$, and
$a^{(l)}\in\mathcal{A}_{D_{i_l}}$ is the action the expert deems
optimal under $\mathbf{x}^{(l)}$. Components marked as
``irrelevant'' contribute no constraint and the rule must hold for
every value they may take.
\end{definition}

Given such a rule base we measure the consistency of a parameter
vector with the expert by a simple accuracy:

\begin{definition}[Fitness function]
\label{def:fitness}
Define the vector of expected utilities of decision $D_{i_l}$ under
$\boldsymbol{\Theta}$ and evidence $\mathbf{x}^{(l)}$ as
\[
   \mathbf{u}^{(l)}(\boldsymbol{\Theta}) =
   \bigl(\mathrm{EU}(a \mid \boldsymbol{\Theta}, \mathbf{x}^{(l)})\bigr)_{a\in\mathcal{A}_{D_{i_l}}}.
\]
The generic loss minimised by the EDA is then
\[
   \mathcal{L}(\boldsymbol{\Theta}) =
   \sum_{(\mathbf{x}^{(l)},\,D_{i_l},\,a^{(l)})\in\mathcal{R}}
   \ell\!\left(\mathbf{u}^{(l)}(\boldsymbol{\Theta}),\; a^{(l)}\right),
\]
where $\ell$ is a per-rule loss that compares the expected-utility
vector of the target decision with the action prescribed by the
expert. Different choices of $\ell$ recover the fitness variants used
in this work:
\[
\begin{aligned}
   \ell_{\mathrm{binary}}(\mathbf{u}, a)  &= \mathbb{1}\!\left[\textstyle\arg\max_{a'} u_{a'} \neq a\right], \\[4pt]
   \ell_{\mathrm{regret}}(\mathbf{u}, a)  &= \max_{a'} u_{a'} - u_a, \\[4pt]
   \ell_{\mathrm{margin}}(\mathbf{u}, a)  &= \max\!\bigl(0,\; \max_{a'\neq a} u_{a'} + m - u_a\bigr), \\[4pt]
   \ell_{\mathrm{softmax}}(\mathbf{u}, a) &= -\log \operatorname{softmax}(\mathbf{u})_a, \\[4pt]
   \ell_{\mathrm{entropy}}(\mathbf{u}, a) &= -\log \operatorname{softmax}(\mathbf{u})_a + \alpha\, H\!\bigl(\operatorname{softmax}(\mathbf{u})\bigr),
\end{aligned}
\]
where $m$ is the margin, $\alpha$ the entropy-regularisation weight, and
$H(\mathbf{p}) = -\sum_{a'} p_{a'}\log p_{a'}$ the Shannon entropy of the
softmax distribution over actions.

\end{definition}

The parameter-recovery problem is then the constrained search
\[
   \boldsymbol{\Theta}^{\star} =
   \arg\min_{\boldsymbol{\Theta}} \mathcal{L}(\boldsymbol{\Theta}),
\]
with $\boldsymbol{\Theta}$ ranging over the product of all valid
CPTs and utility tables. The next proposition justifies the use of
a derivative-free metaheuristic instead of a gradient-based
optimiser.

\begin{proposition}[Irregularity of the objective]
\label{prop:nonsmooth}
The objective $\mathcal{L}$ of Definition~\ref{def:fitness} is
non-convex with respect to $\boldsymbol{\Theta}$ for every per-rule
loss $\ell$ considered in this work. For the $0$--$1$ (\textsc{binary})
loss it is, in addition, piecewise constant and discontinuous; for the
\textsc{regret} and \textsc{margin} losses it is continuous but
non-differentiable on the set where the optimal action of a decision
switches; for \textsc{softmax} and \textsc{entropy} it is piecewise
smooth.
\end{proposition}

\begin{proof}[Sketch]
\textbf{\emph{Non-convexity.}} For a chain of two chance nodes the expected
utility contains a term proportional to the product of two parameters,
$p\,q$. The map $(p,q)\mapsto p\,q$ is not convex: at $(1,0)$ and
$(0,1)$ it equals $0$, yet at the midpoint $(\tfrac12,\tfrac12)$ it
equals $\tfrac14>0$, violating the convexity inequality. Since
$\mathbf{u}^{(l)}(\boldsymbol{\Theta})$ generically involves such
products (a node multiplies the probabilities of its parents) and a
normalisation by the probability of the evidence, $\mathcal{L}$ is
non-convex.

\textbf{\emph{Regularity.}} Within any region of $\boldsymbol{\Theta}$ where the
optimal action of every decision is fixed, $\mathbf{u}^{(l)}$ is a
rational function of the parameters and is therefore smooth. At the
boundaries between regions the maximising action switches: the value of the optimal policy is non-differentiable there. Composing with the per-rule loss,
the $0$--$1$ loss jumps discontinuously at those boundaries, the
\textsc{regret} and \textsc{margin} losses stay continuous but
non-differentiable, and the \textsc{softmax} and \textsc{entropy} losses
are smooth away from them.
\end{proof}

\noindent A gradient of $\mathcal{L}$ thus exists almost everywhere
(except for the discontinuous \textsc{binary} loss) and could in
principle be obtained by finite differences. We nevertheless adopt a
population-based, derivative-free strategy for three reasons: the
landscape is strongly non-convex and highly degenerate (many parameter vectors reproduce the same rules) so local gradient information does not
lead to the global optimum; the evaluation engine returns only utility
values, so each numerical gradient would cost $O(d)$ additional
evaluations in a space of dimension $d=|\boldsymbol{\Theta}|$; and for the
\textsc{binary} loss the gradient is identically zero almost everywhere.


% ==============================================================
\section{Estimation of Distribution Algorithms}
% ==============================================================

Estimation of Distribution Algorithms
(EDAs)~\cite{larranaga2002eda,larranaga2024edas} are a class of
evolutionary metaheuristics in which the classical
crossover-and-mutation operators are replaced by the
\emph{construction and sampling of an explicit probabilistic model}
of the search space. At every generation a probability distribution
is fitted to the best individuals of the current population and used
to generate the next generation. Algorithm~\ref{alg:eda} states the
generic scheme.
\newpage

\begin{algorithm}[Generic EDA]\hfill
\label{alg:eda}
\begin{algorithmic}[1]
\State Initialise a population $\mathcal{P}_0$ of $N$ individuals from a prior on the search space.
\State $t \leftarrow 0$
\Repeat
    \State Evaluate $f$ on every $\boldsymbol{\Theta}\in\mathcal{P}_t$.
    \State $\mathcal{P}_t^{\text{sel}} \leftarrow$ top-$S$ individuals of $\mathcal{P}_t$.
    \State Estimate a probabilistic model $\hat{p}_t$ based on
           $\mathcal{P}_t^{\text{sel}}$.
    \State Sample $N$ new individuals from $\hat{p}_t$ to form
           $\mathcal{P}_{t+1}$.
    \State $t \leftarrow t+1$
\Until{stopping criterion is met}
\State \Return best individual found.
\end{algorithmic}
\end{algorithm}

The defining choice of an EDA is the family used for $\hat{p}_t$:
richer families capture more structure of the search space at the
expense of a higher estimation cost.

\subsection{Variants in this work}

Because the parameter vector of an ID lives in a continuous space
(every CPT entry and every utility value is a real number), only
continuous EDAs are relevant. The four variants used in this work,
all available through \textsc{EDAspy}~\cite{soloviev2024edaspy}, span the spectrum from a fully factorised parametric model to a
full-covariance multivariate Gaussian, with a structured intermediate
and a non-parametric alternative on the univariate side. We denote an individual by $\boldsymbol{\phi}=(\phi_1,\ldots,\phi_d)$
and the size of the selected elite by
$S=|\mathcal{P}_t^{\text{sel}}|$.

\paragraph*{UMDA\textsubscript{c} --- Univariate Marginal Distribution
Algorithm.} Each dimension is modelled by an independent
Gaussian~\cite{larranaga2002eda}:
\[
    \hat{p}_t(\boldsymbol{\phi})
    = \prod_{j=1}^{d}
    \mathcal{N}\!\bigl(\phi_j;\,
    \hat{\mu}_t^{j},\,(\hat{\sigma}_t^{j})^{2}\bigr),
\]
whose mean and variance are the sample moments of the selected individuals,
\[
    \hat{\mu}_t^{j}
    = \frac{1}{S}\!\sum_{\boldsymbol{\phi}\in\mathcal{P}_t^{\text{sel}}}\!\phi_j,
    \qquad
    (\hat{\sigma}_t^{j})^{2}
    = \frac{1}{S}\!\sum_{\boldsymbol{\phi}\in\mathcal{P}_t^{\text{sel}}}\!
    (\phi_j - \hat{\mu}_t^{j})^{2}.
\]
The model is the cheapest of the family ($O(d\cdot S)$ per generation,
both for fitting and for sampling) but its independence assumption
means it cannot capture any dependency between the parameters of the
diagram: the search effectively assumes that good values for one CPT
entry give no information about good values for another.
\paragraph*{EMNA --- Estimation of Multivariate Normal Algorithm.}
The opposite end of the parametric spectrum: the whole population is
modelled by a single multivariate Gaussian~\cite{larranaga2002eda},
\[
    \hat{p}_t(\boldsymbol{\phi}) =
    \mathcal{N}\!\bigl(\boldsymbol{\phi};\,
    \hat{\boldsymbol{\mu}}_t,\,
    \hat{\boldsymbol{\Sigma}}_t\bigr),
\]
with sample mean and \emph{full covariance},
\[
    \hat{\boldsymbol{\mu}}_t
    = \frac{1}{S}\!\sum_{\boldsymbol{\phi}\in\mathcal{P}_t^{\text{sel}}}\!\boldsymbol{\phi},
    \qquad
    \hat{\boldsymbol{\Sigma}}_t
    = \frac{1}{S}\!\sum_{\boldsymbol{\phi}\in\mathcal{P}_t^{\text{sel}}}\!
    (\boldsymbol{\phi} - \hat{\boldsymbol{\mu}}_t)
    (\boldsymbol{\phi} - \hat{\boldsymbol{\mu}}_t)^{\!\top}.
\]
EMNA captures every pairwise linear dependency between parameters and
is therefore strictly more expressive than UMDA. The price is
double: the cost grows to $O(d^{2}\cdot S)$ for the covariance estimate,
plus a $O(d^{3})$ Cholesky factorisation of $\hat{\Sigma}_t$ used to sample (and the
truncated population must be large enough, $S>d$
to keep $\hat{\boldsymbol{\Sigma}}_t$
non-singular)~\cite{larranaga2002eda}.

\paragraph*{EGNA --- Estimation of Gaussian Networks Algorithm.}
A middle ground between UMDA and EMNA: a multivariate Gaussian whose
factorisation is governed by a Bayesian network $\mathcal{G}$
\emph{learnt from the elite at every
generation}~\cite{larranaga2002eda}. Each variable conditions only
on its parents in $\mathcal{G}$,
\[
    \hat{p}_t(\boldsymbol{\phi}) = \prod_{j=1}^{d}
    \mathcal{N}\!\Bigl(\phi_j;\;
    \beta_{j,0} + \!\!\!\!\sum_{k\,\in\,\mathrm{pa}_{\mathcal{G}}(j)}\!\!\!\!
    \beta_{j,k}\,\phi_k,\;
    \sigma_j^{2}\Bigr),
\]
i.e.\ a linear regression of $\phi_j$ on its graph parents with
Gaussian residuals. The structure $\mathcal{G}$ is chosen with a
score (typically BIC) and a greedy search over
DAGs~\cite{larranaga2002eda}. EGNA recovers UMDA when $\mathcal{G}$
has no edges and approaches EMNA when $\mathcal{G}$ becomes a
complete DAG; in practice the learnt graph is sparse, which keeps
the cost reduced and makes the model robust if 
$S$ is not large enough for EMNA to estimate a full covariance
reliably.

\paragraph*{KEDA (Univariate) --- Kernel-density Estimation of
Distribution Algorithm.} A non-parametric alternative on the
univariate side: each marginal is fitted with a kernel density
estimator from the elite~\cite{soloviev2024edaspy},
\[
    \hat{p}_t(\boldsymbol{\phi}) = \prod_{j=1}^{d}
    \hat{p}_t^{\,j}(\phi_j),
    \qquad
    \hat{p}_t^{\,j}(\phi_j)
    = \frac{1}{S\,h_j}\!\sum_{\boldsymbol{\phi}^{(s)}\in\mathcal{P}_t^{\text{sel}}}\!\!
    K\!\Bigl(\frac{\phi_j - \phi_j^{(s)}}{h_j}\Bigr),
\]
where $K$ is a kernel (typically Gaussian) and $h_j$ is a bandwidth
chosen, for instance, by Silverman's rule. Like UMDA it ignores
cross-parameter dependencies, but it drops the Gaussian assumption
per dimension and can therefore fit skewed or multi-modal marginals
(a setting in which a single Gaussian would smear over distinct
modes). Its cost is again $O(d\cdot S)$, with the bandwidth as an
additional implicit hyperparameter that trades smoothness against
detail.

\medskip

Together the four variants span the spectrum from \emph{fully
factored} (UMDA\textsubscript{c}, KEDA) to \emph{fully joint} (EMNA),
with EGNA as a structured intermediate, and from \emph{parametric}
(UMDA\textsubscript{c}, EMNA, EGNA) to \emph{non-parametric} (KEDA)
on the per-dimension side. Which combination suits the
parameter-recovery problem of this work is far from obvious a
priori.

\subsection{Why the algorithm converges}

The reason an EDA makes progress is intuitive and needs no heavy
machinery. At each generation it keeps only the best fraction of the
population — a \emph{truncation} selection that retains the top
$\mu\in(0,1)$ of the individuals — and refits the probabilistic model
$\hat{p}_t$ to that elite. Because the selected individuals
concentrate in the better-scoring regions of the search space, the
model fitted to them places its mass there, so the next population,
sampled from $\hat{p}_t$, is on average better than the previous one.

Repeating the loop has two coupled effects. The model \emph{drifts}
towards regions of lower cost, generation after generation, while its
variance \emph{contracts} as the surviving individuals become more
alike. The drift is what improves the solutions; the contraction is
what makes the search eventually settle on a concrete configuration.
When the variance has shrunk to the point where new samples barely
differ from one another, the population has converged.

The selection ratio $\mu$ balances these two effects. A small $\mu$
keeps only a handful of top individuals: the model contracts quickly
and the search converges fast, but it may lock onto the first good
region it finds. A large $\mu$ keeps almost the whole population, so
little information is extracted from the cost function and the search
behaves almost like random sampling. A useful value lies in between
and is tuned empirically in Chapter~4.

This argument is deliberately informal. It explains why the
distribution concentrates on good regions, not that the region found
is the global optimum — on the non-convex, degenerate landscape of
Proposition~\ref{prop:nonsmooth} it generally need not be. What
matters in practice is that the EDA reliably drives the cost down
towards high-quality solutions, which is what the experiments of
Chapter~4 confirm. The only additional requirement specific to our
problem is that every sampled individual be turned into a valid set of
CPTs — each row a probability distribution — which is handled at the
decoding stage described in Chapter~\ref{ch:implementation}.



% % ==============================================================

% \section{Summary}


% We have introduced the formal machinery needed to state the
% parameter-recovery problem: the definition of an Influence Diagram,
% its quantitative parameterisation and the notion of an optimal
% policy obtained by the Shachter algorithm. We have then formalised
% the expert rule base, the fitness function and shown
% (Proposition~\ref{prop:nonsmooth}) that the resulting optimisation
% problem is non-convex and non-smooth — hence the choice of a
% derivative-free method. Finally, we have presented the EDA family
% with the variants used in this work, an intuitive account of why the
% search converges, and a comparison with classical genetic algorithms.
% The next chapter turns these theoretical pieces into an executable
% implementation.